% ---------------------------------------------------------------------------
% Zusammenfassung
% Wortgleich uebernommen aus der Confluence-Seite
% "Zusammenfassung / Abstract (Entwurf)" (Space BEH, Seiten-ID 101416961),
% Abschnitt "Zusammenfassung", Stand 04.10.2026.
% Auf Wunsch der Verfasser nur die deutsche Fassung; das englische Abstract
% der Confluence-Seite erscheint nicht in der Arbeit.
% ---------------------------------------------------------------------------
\addchap{Zusammenfassung}

An der Hochschule Bremen gibt es keinen Ort, an dem abgeschlossene studentische Projekte dauerhaft und öffentlich einsehbar sind. Die vorliegende Arbeit untersucht, welche funktionalen und gestalterischen Anforderungen ein hochschulisches E-Portfolio-Portal erfüllen muss, um solche Projekte strukturiert zu präsentieren, auffindbar zu machen und von den beteiligten Nutzergruppen akzeptiert zu werden. Aus der Literatur zu E-Portfolios, dem Technologieakzeptanzmodell und den Usability-Heuristiken nach Nielsen sowie aus einer heuristischen Analyse der Referenzplattform der HfG Schwäbisch Gmünd wurden 36 funktionale und 24 nicht-funktionale Anforderungen abgeleitet. Auf dieser Grundlage entstand ein Portal mit blockbasiertem Inhaltsmodell und einem redaktionellen Freigabeweg aus fünf Zuständen, dessen Regeln der Server durchsetzt. Geprüft wurde es mit lautem Denken in elf Sitzungen, einer heuristischen Evaluation durch zwei Begutachtende, der System Usability Scale (SUS) und einem Fragebogen nach dem Technologieakzeptanzmodell.

Der SUS-Mittelwert liegt mit 87,3 Punkten über dem vorab gesetzten Zielwert von 68, und alle elf Personen bewerten Nützlichkeit und Benutzerfreundlichkeit über dem Skalenmittelpunkt. Dem stehen 30 Usability-Probleme gegenüber, sieben davon schwerwiegend. Von den funktionalen Anforderungen sind 30 erfüllt und sechs teilweise, von den nicht-funktionalen zehn erfüllt, elf teilweise und drei nicht geprüft. Auffindbarkeit fördern verbindbare, einzeln rücknehmbare Ordnungsmerkmale und ein Zugang über die Person; Voraussetzung ist, dass Zustände und Handlungen beschriftet sichtbar sind und dasselbe überall gleich heißt. Das Inhaltsmodell lässt sich mit dem Redaktionsworkflow verbinden, indem die Bearbeitbarkeit an den Zustand des Projekts gebunden und auf dem Server durchgesetzt wird. Die Ergebnisse beruhen auf elf Personen und einer einzigen Testrunde.
